\begin{abstract}
We investigate whether small language models (LLMs) can learn to compute matrix inverses through supervised fine-tuning on synthetic data---a controlled test of whether LLMs can internalize mathematical patterns that are computationally non-trivial for humans to reproduce from memory. Matrix inversion is an ideal probe: it has a verifiable ground truth ($AA^{-1} = I$), requires multi-step algebraic reasoning, and generates output that can look superficially plausible while remaining mathematically wrong.

Using LoRA (Low-Rank Adaptation, rank $r = 16$, $\alpha = 32$) to fine-tune Llama-3.2-3B on 800 supervised ($A$, $A^{-1}$) training examples covering $2 \times 2$ and $3 \times 3$ matrices, we evaluate both the untrained baseline and the fine-tuned model on 100 held-out examples using a multi-dimensional diagnostic suite that distinguishes three levels of success: (1) syntactic well-formedness, (2) coarse numeric resemblance to the true inverse, and (3) exact functional correctness via residual-based metrics.

Our main findings are:
\begin{itemize}
    \item \textbf{Baseline}: 100\% of outputs are syntactically parseable and correctly shaped, but 0\% satisfy the inverse constraint. All outputs are classified \texttt{NUMERICALLY\_WRONG}. Average residuals are $\|AX-I\|_F = 1.84$ and $\|XA-I\|_F = 1.85$. The model behaves essentially as a format imitator with no inverse-like numeric content (determinant ratio 0.045, sign agreement 0.503).
    \item \textbf{Fine-tuned}: Fine-tuning on 800 steps improves several coarse numeric diagnostics---determinant ratio rises from 0.045 to 0.639, sign agreement from 0.503 to 0.795, and relative Frobenius error to the gold inverse falls from 1.030 to 0.727---indicating partial internalization of inverse-like numeric structure. However, true inverse accuracy remains 0\% under all residual tolerances, and fine-tuning introduces severe generation instability: only 60\% of outputs are parseable, with 22\% entering repeated-block loops and 18\% producing malformed shapes.
\end{itemize}

The dominant conclusion is that LoRA fine-tuning learned \emph{partial inverse-like numeric structure} but could not bridge the gap to functional correctness, while simultaneously degrading generation reliability. This work provides a controlled diagnostic on the limits of supervised fine-tuning for precise arithmetic computation in small LLMs.
\end{abstract}
